\section{A15.1b1: horizon-open \texorpdfstring{$S = T$}{S=T} endpoint injectivity}
\label{sec:p12-b1}

\begin{theorem}[Horizon-open endpoint injectivity]\label{thm:p12-b1}
Let $2a < T_0 < c$, $0 < R < T$, $S = T$.  Then
\(\ker L_{R,T,T_0}^{\{a,b,2a\}} = \{0\}\).
\end{theorem}

\begin{proof}
Let $h \in L^2(R, T)$ with $L_{T_0} h = 0$.

\medskip
\noindent\emph{Step 1 (lower-half kill from A14.3a persistence).}
For $u = a + x$ with $0 < x < a$, the source E-equation of A14.3a
holds unchanged.  The new upper branch $g(u + 2a) = g(3a + x)$ is
deleted by $P_{T_0}$ because $3a + x > T_0$ (using $T_0 < c < 3a$).
The remaining new forward branch $g(u - 2a) = g(x - a)$ is handled by
oddness $g(x - a) = -g(a - x)$, already inside the A14.3a $q$-reflection
term.  Because the support of $h$ ends at $T$, the fold branch
$g(u + a) = g(2a + x) = g(T + x)$ vanishes for $x < 0$ (trivially) and,
more importantly, for $u > 0$ has $g(2a + x) = 0$ because we assumed
$S = T$ and $h$ vanishes on $(T, T_0)$ by the support hypothesis.

Consequently, the repaired A14.3a lower-half kill applies verbatim and
yields
\[
h = 0 \quad \text{a.e. on } (R, a).
\tag{S1}
\]

\medskip
\noindent\emph{Step 2 (open source horizon activates the fold).}
For $u = T + t$ with $0 < t < \varepsilon = T_0 - T$, all three forward
branches
\[
g(u - a) = g(a + t),\quad
g(u - b) = g(e + t),\quad
g(u - 2a) = g(t)
\]
have arguments above the support of $h$ up to the fold contributions
below.  Applying the E-equation shape:
\[
q\,g(t) + r\,g(e + t) + p\,g(a + t) = 0.
\tag{1}
\]
Since $t < \varepsilon < e < a$, we have $t \in (0, e) \subset (0, a)$
and, if $t > R$, then $t \in (R, a)$; combined with $e + t \in (e, 2e)
\subset (e, a)$ and $t > R$ or $t < R$: both $t$ and $e + t$ live in the
already-killed lower half or (for $t < R$) outside the support.  In
either case,
\[
g(t) = 0, \qquad g(e + t) = 0
\]
by (S1) or by support.  Equation (1) therefore reduces to
$p\, g(a + t) = 0$, hence
\[
h(a + t) = 0 \quad\text{for a.e. } 0 < t < \varepsilon.
\tag{S2}
\]
This produces an open null-strip $(a, a + \varepsilon)$ in the
upper-half support.

\medskip
\noindent\emph{Step 3 (transport by the high reflection).}
Define $l(z) := h(T - z) = h(2a - z)$ for $z \in (0, T - R)$.  The old
$S < T$ upper support gap gave the seed $l(z) = 0$ on $(0, T - S)$; that
gap has vanished at $S = T$.  The new null-strip (S2) gives instead the
\emph{right-sided} gap
\[
l(z) = 0 \quad \text{for } z \in (a - \varepsilon, a).
\tag{3}
\]
The A14.3a high-reflection identity (valid for $d < z < a$) reads
\[
q\, l(z) + p\, l(a - z) = 0.
\tag{4}
\]
Because $\varepsilon < e$ implies $a - \varepsilon > d$, the whole gap
(3) lies in the region where (4) applies.  Substituting (3) into (4)
gives $p\, l(a - z) = 0$, hence $l(a - z) = 0$, i.e., under
$w := a - z$,
\[
l(w) = 0 \quad \text{for } 0 < w < \varepsilon.
\tag{S3}
\]

\medskip
\noindent\emph{Step 4 (upper-circle restart via UC2).}
The upper-circle vector
\[
W(t) := \bigl(l(t),\; l(e - t)\bigr),\qquad 0 < t < e,
\]
has $l(t) = 0$ on the nonempty open subinterval
$(0, \varepsilon) \subset (0, e)$ by (S3).  The matrices, irrational
rotations and non-invariance conditions of the A14.3a upper-circle
step are unchanged from the frozen A14.3a apparatus.  In particular,
the UC2 return-count repair of 2026-08-21 applies to the same
cocycle, delivering
\[
W = 0 \quad \text{a.e. on } (0, e),
\]
whence $h = 0$ on $(a, T)$.  Together with (S1) this yields $h = 0$
a.e. on $(R, T)$.
\end{proof}

\begin{remark}[Structural content]\label{rem:p12-b1-structural}
Theorem~\ref{thm:p12-b1} is not a mere boundary extension of A14.3a.
The A14.3a upper-circle seed was $T - S > 0$; at $S = T$ that seed
vanishes identically.  What replaces it is a source-horizon-generated
null-strip $(a, a + \varepsilon)$ produced by the newly active
$2a$-shift on the fold side, which the pre-existing A14.3a
high-reflection identity transports back to the left seed
$l(w) = 0$ on $(0, \varepsilon)$.  The mechanism
\[
\text{old support gap}\; \leadsto\; \text{new horizon-generated gap}
\]
is genuinely new.
\end{remark}

\begin{remark}[Scope: the exact corner $T_0 = T$ is not decided]
\label{rem:p12-b1-exact-corner}
The proof uses $\varepsilon = T_0 - T > 0$ essentially, both to open
the H-source in (1) and to place the transported gap in the
high-reflection region for (4).  The exact corner $T_0 = T = S$ is
logically separate and remains \(?[O]\).
\end{remark}
